\documentclass{fiche}
\metadonnees{16}{Le vœu}{Bloc 4 — Hypothèse, parole rapportée, traduction}

\begin{document}
\entete

\objectifs{%
\textbf{Langue} : les quatre types de conditionnelles ; \emph{wish} et \emph{if only}.\par
\textbf{Texte} : Oscar Wilde, \emph{The Picture of Dorian Gray} (1891), chapitre
\textsc{ii}.\par
\textbf{Durée} : 1\,h\,30 — exercices 35 min, texte et questions 30 min, version 25 min.}

%% ===================================================================
\exercice[14 min]{Les conditionnelles}

\rappel{\textbf{Rappel.} Le temps de la subordonnée recule d'un cran à chaque degré
d'irréalité : présent, prétérit, pluperfect. \emph{Will} ne se met jamais après \emph{if}.}

\consigne{Mettez les verbes entre parenthèses aux formes qui conviennent.}

\begin{anglais}
\begin{enumerate}[label=\arabic*.,itemsep=2.2mm,leftmargin=*]
  \item If you \emph{(heat)} \trou{heat} water to 100 degrees, it \emph{(boil)} \trou{boils}.
    \emph{(type 0)}
  \item If he \emph{(sell)} \trou{sells} the portrait, Basil \emph{(be)} \trou{will be}
    furious. \emph{(type 1)}
  \item If I \emph{(be)} \trou{were} always young, I \emph{(give)} \trou{would give} my soul
    for it. \emph{(type 2)}
  \item If Dorian \emph{(not meet)} \trou{had not met} Lord Henry, nothing \emph{(happen)}
    \trou{would have happened}. \emph{(type 3)}
  \item If the picture \emph{(can change)} \trou{could change}, I \emph{(be)} \trou{would be}
    always what I am now.
  \item If you \emph{(ask)} \trou{had asked} me yesterday, I \emph{(say)} \trou{would have
    said} no.
  \item What \emph{(you / do)} \trou{would you do} if you \emph{(be)} \trou{were} in his
    place?
  \item If he \emph{(not be)} \trou{had not been} so vain, he \emph{(be)} \trou{would be}
    happy today. \emph{(type mixte)}
  \item I \emph{(help)} \trou{will help} you if you \emph{(want)} \trou{want} me to.
  \item If she \emph{(know)} \trou{had known} the truth, she \emph{(never come)} \trou{would
    never have come}.
  \item Unless you \emph{(hurry)} \trou{hurry}, we \emph{(miss)} \trou{will miss} the train.
  \item If I \emph{(be)} \trou{were} you, I \emph{(not touch)} \trou{would not touch} it.
\end{enumerate}
\end{anglais}

\note[Trois règles et une exception]{Jamais \emph{will} ni \emph{would} après \emph{if} : le
décalage des temps suffit. \emph{Were} à toutes les personnes dans le type 2
(\emph{if I were you}) --- c'est un ancien subjonctif, encore obligatoire dans cette
expression. \emph{Unless} = \emph{if\ldots not} et ne se met donc jamais lui-même au négatif.
L'exception : le type mixte (8), hypothèse passée et conséquence présente, qui combine
pluperfect et conditionnel présent.}

%% ===================================================================
\exercice[10 min]{\emph{Wish} et \emph{if only}}

\rappel{\textbf{Rappel.} Après \emph{wish} et \emph{if only}, le regret recule d'un cran :
prétérit pour le présent, pluperfect pour le passé. \emph{Would} y exprime le reproche.}

\consigne{Complétez.}

\begin{anglais}
\begin{enumerate}[label=\arabic*.,itemsep=2.2mm,leftmargin=*]
  \item I wish I \trou{were} always young. \emph{(je ne le suis pas)}
  \item He wishes he \trou{had never seen} the portrait. \emph{(il l'a vu)}
  \item If only the picture \trou{could} change!
  \item I wish you \trou{would stop} talking about my nerves. \emph{(reproche)}
  \item She wishes she \trou{had} more time.
  \item If only he \trou{had listened} to his friend that evening!
  \item I wish I \trou{had not said} that.
  \item If only I \trou{knew} what he is thinking!
  \item He wishes he \trou{could} paint like Basil.
  \item I wish they \trou{would not stare} at me like that. \emph{(reproche)}
\end{enumerate}
\end{anglais}

\note[Le piège]{\emph{I wish I knew} ne signifie pas \og je souhaite savoir \fg{} mais
\og si seulement je savais \fg. Le prétérit n'y marque pas le passé mais l'irréel : c'est le
même prétérit modal que dans \emph{if I were you}. Et \emph{I wish you would} s'adresse
toujours à quelqu'un dont on attend un changement de comportement : c'est un reproche poli.}

%% ===================================================================
\exercice[11 min]{La jeunesse et le temps}
\consigne{a) Associez chaque mot à sa traduction.}

\begin{multicols}{2}
\begin{enumerate}[label=\arabic*.,itemsep=1.2mm,leftmargin=*]
  \item a wrinkle \hfill \trou{e}
  \item to mar \hfill \trou{h}
  \item a pang \hfill \trou{b}
  \item to quiver \hfill \trou{k}
  \item a lad \hfill \trou{a}
  \item to mock \hfill \trou{g}
  \item to stir \hfill \trou{j}
  \item brevity \hfill \trou{c}
  \item loveliness \hfill \trou{l}
  \item to well up \hfill \trou{d}
  \item to fling \hfill \trou{i}
  \item hard lines \hfill \trou{f}
\end{enumerate}
\columnbreak
\begin{enumerate}[label=\alph*.,itemsep=1.2mm,leftmargin=*]
  \item un garçon, un gamin
  \item un élancement, un pincement
  \item la brièveté
  \item monter (larmes), jaillir
  \item une ride
  \item pas de chance, dur pour
  \item se moquer de, narguer
  \item gâcher, abîmer
  \item jeter, lancer
  \item remuer, émouvoir
  \item frémir
  \item la beauté, la grâce
\end{enumerate}
\end{multicols}

\note[Deux suffixes de noms]{\emph{-ness} transforme un adjectif en nom :
\emph{lovely $\rightarrow$ loveliness}, \emph{coarse $\rightarrow$ coarseness} (fiche 15),
\emph{happy $\rightarrow$ happiness}. \emph{-ity} fait de même à partir d'adjectifs
d'origine latine : \emph{brief $\rightarrow$ brevity}, \emph{cruel $\rightarrow$ cruelty},
\emph{able $\rightarrow$ ability}. Le mot change alors souvent de forme, comme en français.}

\consigne{b) Complétez avec un mot de la liste.}
\begin{anglais}
\begin{enumerate}[label=\arabic*.,itemsep=2mm,leftmargin=*]
  \item Lord Henry's warning about the \trou{brevity} of youth had \trou{stirred} him.
  \item A sharp \trou{pang} of pain struck through him like a knife.
  \item He will like me till I have my first \trou{wrinkle}.
  \item The hot tears \trou{welled} into his eyes and he \trou{flung} himself on the divan.
\end{enumerate}
\end{anglais}

%% ===================================================================
\newpage
\section{Texte}

\noindent\small\textit{Basil Hallward vient d'achever le portrait de Dorian Gray, un jeune
homme d'une grande beauté. Lord Henry Wotton, cynique et brillant, vient de lui expliquer
que la jeunesse ne dure pas.}\normalsize

\begin{texte}
``Don't you like it?'' cried Hallward at last, stung a little by the \gl[a lad]{lad}{un
garçon}'s silence, not understanding what it meant.

``Of course he likes it,'' said Lord Henry. ``Who wouldn't like it? It is one of the greatest
things in modern art. I will give you anything you like to ask for it. I must have it.''

``It is not my property, Harry.''

``Whose property is it?''

``Dorian's, of course,'' answered the painter.

``He is a very lucky fellow.''

``How sad it is!'' murmured Dorian Gray with his eyes still fixed upon his own portrait. ``How
sad it is! I shall grow old, and horrible, and dreadful. But this picture will remain always
young. It will never be older than this particular day of June.\ldots{} If it were only the
other way! If it were I who was to be always young, and the picture that was to grow old! For
that --- for that --- I would give everything! Yes, there is nothing in the whole world I
would not give! I would give my soul for that!''

``You would hardly care for such an arrangement, Basil,'' cried Lord Henry, laughing. ``It
would be rather \gl[hard lines]{hard lines}{dur, pas de chance} on your work.''

``I should object very strongly, Harry,'' said Hallward.

Dorian Gray turned and looked at him. ``I believe you would, Basil. You like your art better
than your friends. I am no more to you than a green bronze figure. Hardly as much, I dare
say.''

The painter stared in amazement. It was so unlike Dorian to speak like that. What had
happened? He seemed quite angry. His face was flushed and his cheeks burning.

``Yes,'' he continued, ``I am less to you than your ivory Hermes or your silver
\gl[a faun]{Faun}{un faune}. You will like them always. How long will you like me? Till I have
my first \gl[a wrinkle]{wrinkle}{une ride}, I suppose. I know, now, that when one loses one's
good looks, whatever they may be, one loses everything. Your picture has taught me that. Lord
Henry Wotton is perfectly right. Youth is the only thing worth having.''

Hallward turned pale and caught his hand. ``Dorian! Dorian!'' he cried, ``don't talk like
that. I have never had such a friend as you, and I shall never have such another. You are not
jealous of material things, are you? --- you who are finer than any of them!''

``I am jealous of everything whose beauty does not die. I am jealous of the portrait you have
painted of me. Why should it keep what I must lose? Every moment that passes takes something
from me and gives something to it. Oh, if it were only the other way! If the picture could
change, and I could be always what I am now! Why did you paint it? It will
\gl[to mock]{mock}{narguer, se moquer de} me some day --- mock me horribly!'' The hot tears
\gl[to well up]{welled}{monter, jaillir} into his eyes; he tore his hand away and,
\gl[to fling, flung]{flinging}{jeter, lancer} himself on the divan, he buried his face in the
cushions, as though he was praying.

``This is your doing, Harry,'' said the painter bitterly.

Lord Henry shrugged his shoulders. ``It is the real Dorian Gray --- that is all.''
\end{texte}
\source{Oscar Wilde, \emph{The Picture of Dorian Gray}, Londres, Ward, Lock \& Co., 1891,
ch.~\textsc{ii}.}

\partiel[15 min]{Questions}
\consigne{\en{Answer in English, in full sentences.}}

\begin{enumerate}[label=\arabic*.,itemsep=2mm,leftmargin=*]
  \item \en{Why does Basil think Dorian does not like the portrait?}
    \lignes{2}
    \corr{Because Dorian says nothing at all. The silence stings him, and he cannot guess
    what it means.}
  \item \en{What exactly does Dorian find sad?}
    \lignes{3}
    \corr{Not the picture, but the comparison: he will grow old, horrible and dreadful, while
    the painted face will stay exactly as it is on this day of June.}
  \item \en{State his wish in your own words, and say what he offers for it.}
    \lignes{3}
    \corr{That the ageing should be the other way round --- himself always young and the
    portrait growing old. He says there is nothing in the world he would not give, and
    offers his soul.}
  \item \en{How does Dorian's attitude to Basil change in the course of the scene?}
    \lignes{3}
    \corr{He begins silent and admiring, then accuses Basil of caring more for his art than
    for his friends, calls himself less to him than a bronze figure, and finally weeps and
    blames the picture. The painter no longer recognises him.}
  \item \en{``This is your doing, Harry.'' Is Basil right?}
    \lignes{3}
    \corr{Lord Henry has just given the speech about the brevity of youth, so he provided the
    idea; but he answers that this is simply the real Dorian Gray. The scene leaves the
    question open, which is the whole subject of the novel.}
  \item \en{Find the sentences in which Dorian speaks of the portrait as if it were alive.}
    \lignes{3}
    \corr{\emph{this picture will remain always young}, \emph{the picture that was to grow
    old}, \emph{everything whose beauty does not die}, \emph{Every moment that passes takes
    something from me and gives something to it}, \emph{It will mock me some day}.}
\end{enumerate}

\note[Les hypothèses du texte]{\emph{If it were only the other way!} et \emph{If the picture
could change, and I could be always what I am now!} : hypothèses irréelles sans principale ---
ce sont des souhaits, équivalents de \emph{I wish\ldots} \emph{Who wouldn't like it?},
\emph{You would hardly care}, \emph{I should object} : conditionnels dont la subordonnée est
sous-entendue (\og si on te le proposait \fg).}

%% ===================================================================
\newpage
\section{Version}
\consigne{Traduisez le passage en français. Il précède immédiatement le texte : Dorian
contemple son portrait sans rien dire.}

\begin{version}
Then had come Lord Henry Wotton with his strange \gl[a panegyric]{panegyric}{un éloge} on
youth, his terrible warning of its \gl[brevity]{brevity}{la brièveté}. That had
\gl[to stir]{stirred}{remuer, émouvoir} him at the time, and now, as he stood gazing at the
shadow of his own \gl[loveliness]{loveliness}{la beauté, la grâce}, the full reality of the
description flashed across him. Yes, there would be a day when his face would be wrinkled and
\gl[wizen]{wizen}{parcheminé, ratatiné}, his eyes dim and colourless, the grace of his figure
broken and deformed. The scarlet would pass away from his lips and the gold steal from his
hair. The life that was to make his soul would \gl[to mar]{mar}{gâcher, abîmer} his body. He
would become dreadful, hideous, and uncouth.

As he thought of it, a sharp \gl[a pang]{pang}{un élancement} of pain struck through him like
a knife and made each delicate fibre of his nature \gl[to quiver]{quiver}{frémir}. His eyes
deepened into \gl[amethyst]{amethyst}{l'améthyste}, and across them came a mist of tears. He
felt as if a hand of ice had been laid upon his heart.
\end{version}
\source{\emph{Ibid.}}

\lignes{20}

\corr{\textbf{Proposition de traduction.} Puis était venu lord Henry Wotton, avec son
étrange éloge de la jeunesse, son terrible avertissement sur sa brièveté. Cela l'avait remué
sur le moment, et maintenant, tandis qu'il se tenait là à contempler l'ombre de sa propre
beauté, toute la réalité de cette description le traversa d'un éclair. Oui, il viendrait un
jour où son visage serait ridé et parcheminé, ses yeux ternes et sans couleur, la grâce de sa
silhouette rompue et déformée. L'écarlate s'effacerait de ses lèvres et l'or déserterait ses
cheveux. La vie, qui devait former son âme, abîmerait son corps. Il deviendrait affreux,
hideux, difforme.

Comme il y songeait, un élancement aigu le transperça comme un couteau et fit frémir chaque
fibre délicate de sa nature. Ses yeux se firent plus sombres, couleur d'améthyste, et un
voile de larmes les recouvrit. Il eut le sentiment qu'une main de glace s'était posée sur son
cœur.}

\note[Choix de traduction 1]{\emph{Then had come Lord Henry Wotton} : inversion après
l'adverbe, courante au pluperfect dans le récit littéraire. Le français peut la garder
(\og Puis était venu lord Henry \fg), ce qui conserve le rythme.}

\note[Choix de traduction 2]{\emph{there would be a day when\ldots} : \emph{would} est ici le
futur vu du passé, non un conditionnel de politesse. Le français emploie le conditionnel
présent, qui sert aux deux --- d'où l'ambiguïté qu'il faut lever par le contexte.}

\note[Choix de traduction 3]{\emph{and the gold steal from his hair} : \emph{steal} dépend
encore de \emph{would}, sous-entendu par économie. Il faut le rétablir en français :
\og et l'or déserterait ses cheveux \fg. Repérer les ellipses de l'auxiliaire est le
principal travail de lecture en version.}

\note[Choix de traduction 4]{\emph{The life that was to make his soul} : \emph{be to} +
infinitif, futur de destin (fiche 6). \og La vie, qui devait former son âme \fg{} --- avec
virgules, car la relative est explicative.}

\note[Choix de traduction 5]{\emph{He felt as if a hand of ice had been laid upon his
heart} : passif + pluperfect après \emph{as if}. Le français préfère l'actif pronominal
(\og une main de glace s'était posée \fg) au passif littéral.}

%% ===================================================================
\partiel{Lexique à mémoriser}
\begin{itemize}[label=--,itemsep=0.8mm,leftmargin=*]\small
  \item \textbf{a lad} — un garçon ; \textit{féminin ancien} a lass
  \item \textbf{a fellow} — un type, un garçon \textit{(fiche 7)}
  \item \textbf{a wrinkle} — une ride ; \textit{verbe} to wrinkle \textit{(fiche 14)}
  \item \textbf{good looks} — la beauté, le physique
  \item \textbf{loveliness} — la grâce ; lovely + -ness
  \item \textbf{brevity} — la brièveté ; brief + -ity
  \item \textbf{a warning} — un avertissement \textit{(fiche 9)}
  \item \textbf{to stir} — remuer, émouvoir
  \item \textbf{to mar} — gâcher ; \textit{proche de} to spoil
  \item \textbf{dim} — terne, faible \textit{(fiche 12)}
  \item \textbf{uncouth} — fruste, difforme \textit{(fiche 12)}
  \item \textbf{a pang} — un élancement ; a pang of remorse
  \item \textbf{to quiver} — frémir
  \item \textbf{to gaze at} — contempler ; \textit{plus long que} to glance at
  \item \textbf{to stare at} — fixer \textit{(fiche 2)}
  \item \textbf{to mock} — narguer ; \textit{nom} mockery
  \item \textbf{to fling} — jeter ; flung, flung
  \item \textbf{to shrug one's shoulders} — hausser les épaules
  \item \textbf{worth} — qui vaut ; worth having\textit{, worth it}
  \item \textbf{one's doing} — l'œuvre de quelqu'un ; this is your doing
\end{itemize}

\end{document}
